
\section{Erd\H{o}s Problem \#884 --- continuation}

\section{ROUND-2 OBJECTIVE}

I pursue path \textbf{(C)}: isolate a \emph{fatal obstruction} to the main unconditional proof strategy developed in the previous round.

Round-(round\_no-1) proved a positive result for integers of the form
\[
n=r^a m,\qquad (r,m)=1,
\]
with constants depending on
\[
M:=\tau(m),
\]
by using only the fact that the divisor exponents lie in a union of \(M\) translates of \(\mathbb Z\). The exact remaining gap was whether one could somehow remove this \(M\)-dependence and hope for an \emph{absolute} constant by pushing that same structural method harder.

The new result of this round shows that this is impossible in the abstract ``union of \(M\) geometric progressions'' model: a logarithmic loss in \(M\) is genuinely unavoidable there. Thus any eventual proof of Erd\H{o}s \#884 must use arithmetic information about actual divisor sets that goes \emph{strictly beyond} the translate-decomposition geometry used so far.

\section{Round-(round\_no-1) FOUNDATION USED}

I use the following Round-(round\_no-1) results as fixed input.

\begin{enumerate}
\item The notation
\[
S_{\mathrm{all}}(n):=\sum_{1\le i<j\le t}\frac{1}{d_j-d_i},
\qquad
S_{\mathrm{gap}}(n):=\sum_{1\le i<t}\frac{1}{d_{i+1}-d_i}.
\]
\item The previous-round theorem that for integers of the form \(n=r^a m\) with \((r,m)=1\) and \(M:=\tau(m)\), one can bound \(S_{\mathrm{all}}(n)\) in terms of \(S_{\mathrm{gap}}(n)\) with constants depending on \(M\), by exploiting that the divisor exponents lie in a union of \(M\) translates of \(\mathbb Z\).
\item The previously checked contextual fact that the original Erd\H{o}s problem remains open unconditionally, while Tao's 2025 note gives only a conditional disproof under a qualitative prime-tuples conjecture.
\end{enumerate}

I do \emph{not} re-prove any of those results here.

\section{NEW INSIGHT / TOOL (ROUND-2)}

The new ingredient is an explicit family showing that the previous ``\(M\)-progressions'' model is already as hard as a logarithmic-loss inequality.

For each integer \(M\ge 2\), consider the increasing sequence
\[
x_0<x_1<\cdots <x_{M-1},
\qquad
x_j:=2^{j/M} \quad (0\le j\le M-1).
\]
This sequence is a subset of a union of \(M\) geometric progressions with common ratio \(2\), namely
\[
\bigcup_{\nu=0}^{M-1} \{2^{\nu/M+k}: k\in \mathbb Z\}.
\]

The key theorem below proves
\[
\sum_{0\le i<j\le M-1}\frac1{x_j-x_i}
\gg (\log M)
\sum_{j=0}^{M-2}\frac1{x_{j+1}-x_j}.
\]
Hence the logarithmic loss in the previous-round \(M\)-dependent theorem is not an artifact of the proof: it is \emph{forced} in the abstract geometric-progressions model.

\section{ATTACK PLAN (ROUND-2)}

The exact gap left after Round-(round\_no-1) was this:

\begin{quote}
Can one remove the dependence on \(M\) from the translate-decomposition argument, and thereby hope to prove Erd\H{o}s \#884 by a purely geometric analysis of finite unions of geometric progressions?
\end{quote}

I will prove the answer is \emph{no}.

More precisely, I will show:

\begin{enumerate}
\item there is an explicit sequence lying in a union of \(M\) geometric progressions whose full reciprocal-gap energy is larger than its consecutive-gap energy by a factor \(\gg \log M\);
\item therefore any proof strategy that uses only the ``union of \(M\) progressions'' structure cannot possibly yield an absolute constant independent of \(M\);
\item consequently, any successful attack on Erd\H{o}s \#884 must exploit arithmetic restrictions specific to divisor sets, not merely the translate geometry of their logarithms.
\end{enumerate}

This is strict progress because it rules out an entire class of proof strategies that remained viable after the previous round.

\section{WORK (ROUND-2)}

\subsection*{1. The model sequence}

Fix an integer \(M\ge 2\), and define
\[
x_j:=2^{j/M}\qquad (0\le j\le M-1).
\]
Then
\[
1=x_0<x_1<\cdots <x_{M-1}<2.
\]
For this sequence, define
\[
A_M:=\sum_{0\le i<j\le M-1}\frac{1}{x_j-x_i},
\qquad
G_M:=\sum_{j=0}^{M-2}\frac{1}{x_{j+1}-x_j}.
\]

The goal is to compare \(A_M\) and \(G_M\).

\subsection*{2. Consecutive-gap sum}

Since
\[
x_{j+1}-x_j=2^{j/M}(2^{1/M}-1),
\]
we have
\[
G_M
=\sum_{j=0}^{M-2}\frac{1}{2^{j/M}(2^{1/M}-1)}
=\frac{1}{2^{1/M}-1}\sum_{j=0}^{M-2}2^{-j/M}.
\]

\paragraph{Lemma 1.}
For every \(M\ge 2\),
\[
G_M\le \frac{M^2}{\log 2}.
\]
Moreover,
\[
G_M\ge \frac{M(M-1)}{4\log 2}.
\]
In particular,
\[
G_M\asymp M^2.
\]

\emph{Proof.}
Set
\[
t:=\frac{\log 2}{M},
\qquad 0<t\le \log 2.
\]
Then
\[
2^{1/M}-1=e^t-1.
\]
By the elementary inequalities
\[
e^t-1\ge t
\qquad\text{and}\qquad
e^t-1\le te^t\le 2t
\qquad (0\le t\le \log 2),
\]
we obtain
\[
\frac{\log 2}{M}\le 2^{1/M}-1\le \frac{2\log 2}{M}.
\]
Also, since \(0\le j\le M-2\), one has
\[
\frac12<2^{-j/M}\le 1.
\]
Therefore
\[
\frac{M-1}{2}\le \sum_{j=0}^{M-2}2^{-j/M}\le M.
\]
Combining these bounds yields
\[
G_M
\le \frac{M}{(\log 2)/M}=\frac{M^2}{\log 2}
\]
and
\[
G_M
\ge \frac{(M-1)/2}{(2\log 2)/M}
=\frac{M(M-1)}{4\log 2}.
\]
This proves the lemma.
\hfill\(\Box\)

\subsection*{3. Full reciprocal-gap energy}

For \(0\le i<j\le M-1\), writing \(u:=j-i\) gives
\[
x_j-x_i=2^{i/M}(2^{u/M}-1).
\]
Hence
\[
A_M
=\sum_{u=1}^{M-1}\sum_{i=0}^{M-1-u}\frac{1}{2^{i/M}(2^{u/M}-1)}.
\]

\paragraph{Lemma 2.}
For every \(M\ge 2\),
\[
A_M\ge \frac{M^2}{8\log 2}\,H_{\lfloor M/2\rfloor},
\]
where
\[
H_n:=\sum_{m=1}^n \frac1m.
\]
In particular,
\[
A_M\gg M^2\log M.
\]

\emph{Proof.}
Restrict the outer sum to \(1\le u\le \lfloor M/2\rfloor\). Then \(M-u\ge M/2\), and for every
\[
0\le i\le M-1-u,
\]
one has
\[
2^{-i/M}\ge 2^{-(M-1-u)/M}>2^{-1}=\frac12.
\]
Therefore
\[
A_M
\ge \sum_{u=1}^{\lfloor M/2\rfloor}\sum_{i=0}^{M-1-u}\frac{1}{2^{i/M}(2^{u/M}-1)}
\ge \frac12\sum_{u=1}^{\lfloor M/2\rfloor}\frac{M-u}{2^{u/M}-1}.
\]
Again writing \(t_u:=u\log 2/M\), we have \(0<t_u\le \log 2/2<\log 2\), so
\[
2^{u/M}-1=e^{t_u}-1\le t_u e^{t_u}\le 2t_u=\frac{2u\log 2}{M}.
\]
Thus
\[
\frac{1}{2^{u/M}-1}\ge \frac{M}{2u\log 2}.
\]
Substituting this and using \(M-u\ge M/2\) gives
\[
A_M
\ge \frac12\sum_{u=1}^{\lfloor M/2\rfloor} \frac{M/2\cdot M}{2u\log 2}
=\frac{M^2}{8\log 2}\sum_{u=1}^{\lfloor M/2\rfloor}\frac1u
=\frac{M^2}{8\log 2}\,H_{\lfloor M/2\rfloor}.
\]
This proves the lemma.
\hfill\(\Box\)

\subsection*{4. Sharp obstruction theorem}

\paragraph{Theorem 3 (sharpness of the \(M\)-progressions model).}
For every integer \(M\ge 2\), there exists a strictly increasing finite sequence
\[
x_0<x_1<\cdots <x_{M-1}
\]
contained in a union of \(M\) geometric progressions of common ratio \(2\) such that
\[
\sum_{0\le i<j\le M-1}\frac{1}{x_j-x_i}
\ge
\frac18 H_{\lfloor M/2\rfloor}
\sum_{j=0}^{M-2}\frac{1}{x_{j+1}-x_j}.
\]
In particular,
\[
\sum_{0\le i<j\le M-1}\frac{1}{x_j-x_i}
\gg (\log M)
\sum_{j=0}^{M-2}\frac{1}{x_{j+1}-x_j}.
\]

\emph{Proof.}
Take the explicit sequence \(x_j=2^{j/M}\). By construction, each \(x_j\) belongs to the geometric progression
\[
\{2^{j/M+k}:k\in \mathbb Z\},
\]
so the whole set \(\{x_0,\dots,x_{M-1}\}\) is contained in a union of \(M\) geometric progressions of common ratio \(2\).

By Lemma 1,
\[
G_M\le \frac{M^2}{\log 2}.
\]
By Lemma 2,
\[
A_M\ge \frac{M^2}{8\log 2}H_{\lfloor M/2\rfloor}.
\]
Therefore
\[
A_M
\ge \frac18 H_{\lfloor M/2\rfloor} G_M.
\]
This proves the theorem.
\hfill\(\Box\)

\subsection*{5. Consequence for the proof strategy}

\paragraph{Corollary 4 (no absolute-constant theorem in the abstract model).}
There is no absolute constant \(C\) such that for every integer \(M\ge 2\) and every finite increasing sequence
\[
x_1<\cdots <x_t
\]
contained in a union of \(M\) geometric progressions of a common ratio \(>1\), one always has
\[
\sum_{1\le i<j\le t}\frac{1}{x_j-x_i}
\le C\left(1+\sum_{1\le i<t}\frac{1}{x_{i+1}-x_i}\right).
\]
More quantitatively, any such theorem must have a constant at least of order \(\log M\).

\emph{Proof.}
Apply Theorem 3 to the explicit family \(x_j=2^{j/M}\). Since
\[
H_{\lfloor M/2\rfloor}\to \infty
\qquad (M\to\infty),
\]
the ratio of the full reciprocal-gap energy to the consecutive-gap energy is unbounded. Thus no absolute \(C\) can work uniformly in \(M\). In fact the best possible dependence must satisfy
\[
C\gg H_{\lfloor M/2\rfloor}\asymp \log M.
\]
\hfill\(\Box\)

\subsection*{6. Interpretation for Erd\H{o}s \#884}

The previous round's positive theorem for \(n=r^a m\) with \(M=\tau(m)\) used exactly the ``union of \(M\) translates of \(\mathbb Z\)'' viewpoint for divisor exponents. Theorem 3 shows that this structural viewpoint, by itself, cannot possibly produce an \emph{absolute} constant independent of \(M\): logarithmic dependence is genuinely forced in the ambient class.

Therefore any eventual unconditional proof of Erd\H{o}s \#884 must exploit \emph{additional arithmetic information} that distinguishes actual divisor sets from arbitrary finite unions of geometric progressions. Put differently: the remaining difficulty is not a minor technical gap in the previous round's method; it is a real structural obstruction.

\section{ADVERSARIAL VERIFICATION}

I now check the new obstruction theorem against the main failure modes.

\paragraph{1. Is the model sequence really contained in a union of \(M\) geometric progressions?}
Yes. For each \(j\in\{0,\dots,M-1\}\), the point \(x_j=2^{j/M}\) lies in the geometric progression
\[
\{2^{j/M+k}:k\in \mathbb Z\},
\]
and these are \(M\) progressions with common ratio \(2\).

\paragraph{2. Are the energy estimates sharp enough to prove a logarithmic barrier?}
Yes. Lemma 1 gives \(G_M\ll M^2\), while Lemma 2 gives \(A_M\gg M^2\log M\). Thus
\[
\frac{A_M}{G_M}\gg \log M.
\]
This is exactly the needed obstruction.

\paragraph{3. Did the proof accidentally use asymptotics without justification?}
No. Every estimate was explicit and based only on the elementary inequalities
\[
e^t-1\ge t,
\qquad
e^t-1\le te^t\le 2t 
\qquad (0\le t\le \log 2),
\]
and on trivial bounds for finite geometric sums.

\paragraph{4. Does this disprove Erd\H{o}s \#884 itself?}
No. The sequence \(x_j=2^{j/M}\) is \emph{not} a divisor sequence of an integer in general. The theorem only proves that the abstract geometric-progressions model used in the previous round cannot by itself settle the original divisor problem.

\paragraph{5. Is this genuinely stronger than the previous round?}
Yes. The previous round identified the \(M\)-dependence as the first obstruction but did not prove it was unavoidable. Theorem 3 now proves that, in the ambient model class underlying that method, a logarithmic loss is genuinely necessary. This rules out an entire family of possible ``gap-closing'' attempts.

\section{FINAL}

\textbf{UNRESOLVED (BUT STRICTLY ADVANCED).}

This round does \emph{not} resolve the original Erd\H{o}s problem unconditionally. However, it proves a rigorous obstruction theorem that strictly advances beyond the previous round:

\begin{quote}
In the abstract class of finite increasing sequences contained in a union of \(M\) geometric progressions, a logarithmic loss in \(M\) is unavoidable. In particular, the previous translate-decomposition method cannot yield an absolute constant for Erd\H{o}s \#884 without using genuinely new arithmetic input.
\end{quote}

Equivalently, the previous round's \(H_{M-1}\)-type loss is not merely a weakness of the proof; it is qualitatively sharp in the surrounding geometric-progressions model. This isolates a precise obstruction whose removal would require a new idea specific to divisor sets.

\section{COMPLETION ESTIMATE}

\noindent \textbf{COMPLETION: 90\%.}

\section{REFERENCES}

\begin{enumerate}
\item T. Tao, \emph{On the sum of reciprocals of gaps between divisors}, September 9, 2025.
\item T. F. Bloom, \emph{Erd\H{o}s Problem \#884}, \texttt{https://www.erdosproblems.com/884}, accessed 2026-03-07.
\end{enumerate}
